\documentclass[11pt]{article}
\usepackage[a4paper,margin=27mm]{geometry}
\usepackage{amsmath,amssymb,amsthm,booktabs,array}
\usepackage[hidelinks]{hyperref}
\newtheorem{theorem}{Theorem}
\newtheorem{proposition}[theorem]{Proposition}
\newtheorem{corollary}[theorem]{Corollary}
\theoremstyle{definition}
\newtheorem{definition}[theorem]{Definition}
\newcommand{\ISS}{\mathrm{ISS}}
\newcommand{\dd}{\delta}
\newcommand{\DD}{\Delta}
\title{Graphs covered by vertex and edge identity Seidel switches}
\author{Haobo Ma\\{\small ChronoAI Pte Ltd; The Omega Institute}
\and Wenlin Zhang\\{\small National University of Singapore; The Omega Institute}}
\date{}
\begin{document}
\maketitle
\begin{abstract}
We study finite simple graphs in which every vertex belongs to an identity
Seidel switch consisting of a singleton or the endpoints of an edge.
We prove that every such graph of order $n\geq3$ is connected and has
minimum and maximum degrees satisfying $\delta+\Delta=n$.
Consequently the property is incompatible with complementation and with
self-complementarity. We characterize the bipartite members as the complete
bipartite graphs and give an exact local characterization of the regular
members, together with infinite constructions. An exhaustive enumeration
through order eight gives respectively $1,3,3,11,10,35$ isomorphism classes,
with explicit certificates for every successful singleton and edge switch.
A three-vertex example shows that identity switches need not be closed
under composition.
\end{abstract}

\section{Switches and the cut identity}
Let $G=(V,E)$ be a finite simple graph of order $n$. For $S\subseteq V$,
the Seidel switch $G^S$ complements precisely the adjacencies between $S$
and $V\setminus S$. Following Gervacio~\cite{gervacio}, $S$ is an
\emph{identity Seidel switch} (ISS) if $G^S\cong G$. A singleton ISS is a
vertex-ISS; an edge-ISS is an ISS $\{x,y\}$ with $xy\in E$. We call $G$
\emph{covered} if every vertex belongs to a vertex-ISS or an edge-ISS.
All isomorphisms below may permute the entire vertex set.

Write $e_G(S,V\setminus S)$ for the number of edges crossing the cut.
\begin{proposition}\label{cut}
For every $S\subseteq V$,
\[
 |E(G^S)|-|E(G)|=|S|(n-|S|)-2e_G(S,V\setminus S).
\]
In particular, a vertex-ISS $v$ satisfies $2\deg_G(v)=n-1$, and an
edge-ISS $xy$ satisfies $\deg_G(x)+\deg_G(y)=n$.
\end{proposition}
\begin{proof}
There are $|S|(n-|S|)$ potential cut edges. Switching removes the
$e_G(S,V\setminus S)$ present ones and adds the remaining ones; all other
edges are unchanged. Isomorphism implies equality of edge counts.
For $S=\{v\}$ the cut has size $\deg_G(v)$. For an edge $xy$ and
$S=\{x,y\}$ its size is $\deg_G(x)+\deg_G(y)-2$, so the identity gives
$2(n-2)=2(\deg_G(x)+\deg_G(y)-2)$.
\end{proof}
For a nonadjacent pair the corresponding necessary degree sum is $n-2$.
The adjacency requirement in the definition of edge-ISS is therefore
substantive.

\begin{proposition}\label{singletons}
A nonempty graph has every vertex a vertex-ISS if and only if $n=1$.
\end{proposition}
\begin{proof}
The singleton case is immediate. Otherwise Proposition~\ref{cut} makes
$G$ regular of degree $d=(n-1)/2>0$. Choose an edge $vw$. Switching at
$v$ lowers the degree of $w$ to $d-1$, whereas every graph isomorphic to
$G$ is $d$-regular. This is a contradiction.
\end{proof}
This answers Problem 6.1 of~\cite{gervacio} and explains the need to
allow edge-ISS in the covering question.

\section{Restrictions and two characterizations}
\begin{theorem}\label{extremes}
If $G$ is covered and $n\geq3$, then
\[
 \delta(G)+\Delta(G)=n.
\]
Moreover, $G$ is connected, its complement is not covered, and $G$ is
not self-complementary.
\end{theorem}
\begin{proof}
Put $h=(n-1)/2$, $\delta=\delta(G)$ and $\Delta=\Delta(G)$.
Any vertex whose degree differs from $h$ must belong to an edge-ISS,
and thus has a neighbor of degree $n-\deg_G(v)$.
If $\Delta\leq h$, a vertex of degree less than $h$ would require a
neighbor of degree greater than $h$. Thus all vertices would have degree
$h$. No edge could then be an edge-ISS, since its degree sum would be
$n-1$. Coverage would make all vertices vertex-ISS, contradicting
Proposition~\ref{singletons}. Therefore $\Delta>h$.

A maximum-degree vertex belongs to an edge-ISS, giving
$\delta\leq n-\Delta$. If $\delta\ne h$, a minimum-degree vertex
also belongs to an edge-ISS, giving $n-\delta\leq\Delta$.
If $\delta=h$, then $n$ is odd and the integer $\Delta>h$ satisfies
$\Delta\geq h+1=n-\delta$. In both cases $\delta+\Delta=n$.

If $G$ were disconnected, the component containing a maximum-degree
vertex would have at least $\Delta+1$ vertices and any other component
at least $\delta+1$. Their total would exceed $n$.
Finally,
\[
 \delta(\overline G)+\Delta(\overline G)
 =2(n-1)-(\delta+\Delta)=n-2,
\]
so $\overline G$ cannot be covered. Since coverage is invariant under
isomorphism, $G$ cannot be self-complementary.
\end{proof}

Call adjacent vertices $x,y$ \emph{complementary outside the pair} if
every other vertex is adjacent to exactly one of them.
\begin{theorem}\label{regular}
A regular graph $G$ of order $n\geq3$ is covered if and only if $n$ is
even, its degree is $n/2$, and every vertex has an adjacent vertex
complementary outside the pair. In a regular covered graph these pairs
are exactly the edge-ISS.
\end{theorem}
\begin{proof}
For a regular covered graph Theorem~\ref{extremes} gives $2d=n$.
There are no vertex-ISS, by Proposition~\ref{cut}. Suppose $xy$ is an
edge-ISS. For $w\notin\{x,y\}$, let $c_w$ be the number of its neighbors
in $\{x,y\}$. After switching, its degree is $d+2-2c_w$. The switched
graph must still be $d$-regular, so $c_w=1$.
Conversely, whenever $x,y$ are complementary outside the pair, the
transposition $(x\ y)$, fixing every other vertex, is an isomorphism
$G^{\{x,y\}}\to G$. Such pairs covering all vertices therefore imply
coverage. This sufficiency holds without regularity.
\end{proof}

\begin{theorem}\label{bipartite}
The bipartite covered graphs of order $n\geq3$ are exactly
$K_{a,b}$ with $a,b\geq1$ and $a+b=n$.
\end{theorem}
\begin{proof}
By Theorem~\ref{extremes}, a bipartite covered graph has a bipartition
$A\cup B$ with $1\leq a=|A|\leq b=|B|$. An edge-ISS $xy$ with
$x\in A$, $y\in B$ has degree sum $a+b$, so the bounds
$\deg(x)\leq b$ and $\deg(y)\leq a$ are both attained.
Take any $v\in B$. If it belongs to an edge-ISS, it has degree $a$.
Otherwise it is a vertex-ISS, of degree $(n-1)/2\leq a\leq\lfloor n/2\rfloor$;
this is possible only for odd $n$ and $a=(n-1)/2$, again giving degree
$a$. Every vertex in $B$ is therefore adjacent to all of $A$.
Conversely, every cross-part pair in $K_{a,b}$ is complementary outside
the pair, so every edge is an edge-ISS.
\end{proof}

\section{Constructions and the irregular case}
For each $m\geq2$, partition $2m$ vertices into adjacent pairs
$\{a_i,b_i\}$, $1\leq i\leq m$. Between every two pairs choose either
of the two perfect matchings, independently. Each vertex has degree $m$,
and $a_i,b_i$ are complementary outside their pair. This constructs
regular covered graphs of every even order at least four. Choices include
$K_{m,m}$ and the graph consisting of two copies of $K_m$ joined by a
perfect matching. The latter gives the triangular prism when $m=3$.

More generally, suppose every vertex of a graph has a neighbor
complementary outside the pair. Replace each vertex $u$ by a nonempty
independent set $I_u$, with all edges between $I_u$ and $I_v$ precisely
when $uv$ was an edge. For complementary adjacent $u,v$, every pair
$x\in I_u$, $y\in I_v$ is again complementary outside the pair:
other vertices in these two sets see exactly one endpoint, and all
remaining sets inherit the original property. Thus arbitrary positive
independent-set blow-ups preserve this sufficient condition and give
irregular covered graphs as well.

The sufficient condition is not necessary for irregular graphs. The paw,
with edges $03,12,13,23$, is covered by $\{0,3\}$ and $\{1,2\}$;
the transposition $(0\ 3)$ certifies both switches. The pair $1,2$ has
both a common neighbor and a common nonneighbor, and neither of these
vertices has a complementary neighbor. The house, with edges
$01,03,04,12,23,34$, needs a singleton in its cover: its successful
switches are exactly $\{4\}$, $\{0,1\}$ and $\{2,3\}$.

Theorems~\ref{regular} and~\ref{bipartite} characterize two subclasses
for all orders. A characterization of the remaining irregular,
nonbipartite graphs remains open. The extremal degree identity alone does not characterize coverage.
For example, take $K_{1,4}$ with center $4$ and add the edge $01$.
Its degree sequence is $1,1,2,2,4$, so $\delta+\Delta=5$.
The vertex $0$ has no neighbor of the required degree $3$ for an
edge-ISS. Switching at $0$ changes the degree sequence to $1,2,2,2,3$,
so $0$ is not a vertex-ISS either.

\section{Certified enumeration}
We tested every singleton and every adjacent pair in every simple graph
up to isomorphism for $3\leq n\leq8$. The input consists of the
Read--Wilson graph atlas~\cite{atlas}, as distributed by NetworkX for
orders at most seven, and McKay's catalogue~\cite{mckay} at order eight.
The results are:
\begin{center}
\begin{tabular}{rrrrrr}
\toprule
$n$ & All graphs & Switches tested & Covered & Regular & Bipartite\\
\midrule
3 & 4 & 18 & 1 & 0 & 1\\
4 & 11 & 77 & 3 & 1 & 2\\
5 & 34 & 340 & 3 & 0 & 2\\
6 & 156 & 2106 & 11 & 2 & 3\\
7 & 1044 & 18270 & 10 & 0 & 3\\
8 & 12346 & 271612 & 35 & 3 & 4\\
\bottomrule
\end{tabular}
\end{center}
At orders three through five the members are $K_{1,2}$;
$K_{1,3}$, the paw and $C_4$; and $K_{1,4}$, the house and $K_{2,3}$,
respectively. The regular members of order six are $K_{3,3}$ and the
triangular prism. At orders five and seven, respectively one and six
covered graphs have a vertex belonging to no edge-ISS, so the singleton
option enlarges the class.

For each input graph the computation records a decision for every tested
switch. The cut identity and degree multiset give exact rejection tests;
remaining cases use exhaustive graph isomorphism. Every successful switch
has a stored permutation $p$, with the convention
\[
 uv\in E(G^S)\quad\Longleftrightarrow\quad p(u)p(v)\in E(G).
\]
The accompanying data include graph6 representatives, edge lists and all
such certificates for the covered graphs. A separate standard-library
checker decodes the graphs, toggles the cut and independently searches all
compatible bijections for every tested switch. It also checks every
certificate and verifies pairwise nonisomorphism of the input graphs.
Catalogue completeness through order five is additionally checked by
matching every labelled graph to a representative. For orders six through
eight, completeness uses the cited catalogues; the distinct representative
counts agree with the established counts.

\section{Composition and source distinctions}
Let $G$ have vertices $0,1,2$ and the sole edge $12$. Switching at $1$
gives the sole edge $01$, with certificate $(0\ 2)$; switching at $2$
gives the sole edge $02$, with certificate $(0\ 1)$. Both singletons
are ISS. Their composition is the switch by $\{1,2\}$, which gives
$K_3$ and hence is not an ISS. Thus the ISS family need not be closed
under symmetric difference, even modulo complementary switch sets.
No smaller graph has this failure: for order one all switches are
trivial, and for order two a singleton changes the edge count.

The switch operators themselves do satisfy
$(G^S)^T=G^{S\mathbin{\triangle}T}$. The failure arises because an
isomorphism witnessing $G^S\cong G$ may move the labels of $T$.
The first version of~\cite{gervacio} asserts that all ISS form a group; that assertion is false, and no such assertion is used here.

There is a further distinction relevant to structural characterization.
The paw above contradicts the neighborhood-partition assertion of
Lemma 5.2 in that version. Nor can an isomorphism always be chosen to
exchange the switched endpoints and preserve their complement, as used
in Theorem 5.3 there. For example, take the graph with edges
\[
 01,03,05,12,23,24,25,34,45.
\]
The switch by $\{3,4\}$ has certificate
$p=(3,2,1,5,4,0)$, where the tuple lists the images of $0,\ldots,5$.
On the remaining vertices the induced graph is the cycle $0,1,2,5,0$.
The neighborhood of $3$ there is $\{0,2\}$, a nonedge, whereas the
complement of the neighborhood of $4$ there is $\{0,1\}$, an edge.
No automorphism of this cycle maps the first set to the second.
The cut proof and the unrestricted isomorphism searches above avoid
these assumptions.

\begin{thebibliography}{9}
\bibitem{gervacio} S. V. Gervacio, \emph{On identity Seidel switches},
arXiv:2601.04530v1, 2026. \url{https://arxiv.org/abs/2601.04530v1}.
\bibitem{atlas} R. C. Read and R. J. Wilson, \emph{An Atlas of Graphs},
Oxford University Press, 1998.
\bibitem{mckay} B. D. McKay, \emph{Graphs}, small graph data catalogue,
\url{https://users.cecs.anu.edu.au/~bdm/data/graphs.html}.
\end{thebibliography}
\end{document}
